\documentclass[11pt,a4paper]{article}
\usepackage[margin=27mm,headheight=15pt]{geometry}
\usepackage[T1]{fontenc}
\usepackage[utf8]{inputenc}
\usepackage{libertinus}
\usepackage{amsmath,amssymb,amsthm,mathtools}
\usepackage{microtype,booktabs,array,longtable}
\usepackage{xcolor,enumitem,fancyhdr,titlesec,xurl}
\usepackage{aliascnt}
\usepackage{hyperref}
\usepackage[nameinlink,noabbrev]{cleveref}
\definecolor{ink}{RGB}{25,53,75}
\definecolor{muted}{RGB}{80,88,96}
\hypersetup{colorlinks=true,linkcolor=ink,citecolor=ink,urlcolor=ink,
 pdftitle={Integer Pressure and a Missing Taylor Coefficient},
 pdfauthor={Research draft prepared with ChatGPT for Vladimir Reshetnikov},
 pdfsubject={Analytic phase dependence and the sixth-moment optimum for generalized Thue--Morse products}}
\titleformat{\section}{\Large\bfseries\color{ink}}{\thesection}{.65em}{}
\titleformat{\subsection}{\large\bfseries}{\thesubsection}{.65em}{}
\pagestyle{fancy}\fancyhf{}
\fancyhead[L]{\small Integer pressure and phase optimization}
\fancyhead[R]{\small ProveIt research note}
\fancyfoot[C]{\thepage}
\setlength{\emergencystretch}{2em}
\setlength{\parskip}{.24em}
\setlist{itemsep=.25em,topsep=.4em}
\allowdisplaybreaks
\numberwithin{equation}{section}
\newtheorem{theorem}{Theorem}[section]
\newaliascnt{lemma}{theorem}
\newtheorem{lemma}[lemma]{Lemma}
\aliascntresetthe{lemma}
\newaliascnt{proposition}{theorem}
\newtheorem{proposition}[proposition]{Proposition}
\aliascntresetthe{proposition}
\newaliascnt{corollary}{theorem}
\newtheorem{corollary}[corollary]{Corollary}
\aliascntresetthe{corollary}
\theoremstyle{definition}
\newaliascnt{definition}{theorem}
\newtheorem{definition}[definition]{Definition}
\aliascntresetthe{definition}
\newaliascnt{question}{theorem}
\newtheorem{question}[question]{Question}
\aliascntresetthe{question}
\theoremstyle{remark}
\newaliascnt{remark}{theorem}
\newtheorem{remark}[remark]{Remark}
\aliascntresetthe{remark}
% ed. (2026-09-30): unnumbered environment for ProveIt editorial notes (no counter is shifted).
\newtheorem*{ednote}{Editorial note (ProveIt, 2026-09-30)}
\newcommand{\TT}{\mathbb T}
\newcommand{\RR}{\mathbb R}
\newcommand{\CC}{\mathbb C}
\newcommand{\QQ}{\mathbb Q}
\newcommand{\ZZ}{\mathbb Z}
\newcommand{\NN}{\mathbb N}
\newcommand{\ee}{\mathrm e}
\newcommand{\ii}{\mathrm i}
\newcommand{\dd}{\,\mathrm d}
\newcommand{\e}[1]{\mathrm e^{2\pi\ii #1}}
\newcommand{\Lop}{\mathcal L}
\newcommand{\Em}{\mathcal E_m}
\newcommand{\supp}{\operatorname{supp}}
\newcommand{\spec}{\operatorname{spec}}
\newcommand{\sinc}{\operatorname{sinc}}
\newcommand{\Res}{\operatorname{Res}}
\newcommand{\tr}{\operatorname{tr}}
\newcommand{\rank}{\operatorname{rank}}
\newcommand{\coeff}[2]{[#1^{#2}]}
\newcommand{\repo}{\texttt{ProveIt}}
\newcommand{\code}[1]{\texttt{\detokenize{#1}}}
\newcommand{\eul}[2]{\left\langle\!\begin{matrix}#1\\#2\end{matrix}\!\right\rangle}

\begin{document}
% ed. (2026-09-30): no page anchors on the title page, which is numbered 1 like the following page (removes a duplicate-destination warning of the delivered source).
\hypersetup{pageanchor=false}
\begin{titlepage}
\setlength{\parindent}{0pt}
\thispagestyle{empty}
{\small\sffamily\color{ink} PROVEIT\quad /\quad ANALYSIS AND DYNAMICAL SYSTEMS}\par
\vspace{18mm}
{\Huge\bfseries\color{ink} Integer Pressure and a\\[3mm] Missing Taylor Coefficient\par}
\vspace{7mm}
{\Large Analytic phase dependence and the sixth-moment optimum\\[2mm]
for generalized Thue--Morse products\par}
\vspace{10mm}
{\large Research draft prepared for Vladimir Reshetnikov\par}
{\normalsize Prepared with ChatGPT\quad\textbullet\quad 28 September 2026\par}
\vspace{12mm}
\hrule\vspace{5mm}
\begin{abstract}
For the doubling map and the translated potential
$\psi_c(x)=\log\cos^2\!\pi(x-c)$, we give an explicit finite-dimensional
model of every positive integer pressure $p_m(c)=P(m\psi_c)$.
Its Perron eigenvalue is algebraically simple for every phase, including the
atomic phase $c=0$, and is uniformly separated from the remaining spectrum
when $m$ is fixed. Consequently, every function $c\mapsto p_m(c)$ is
real-analytic. This extends the previously established order-two phase
regularity to all positive integer orders; it does not settle noninteger orders.

The finite model also reveals an all-order cancellation:
$[c^{2m}]p_m(c)=0$. All lower even coefficients are explicit zeta values.
The mechanism is delayed feedback from the second inverse branch, whose
weight first appears in degree $2m$, together with an exact periodized-sinc
Perron eigenfunction at $c=0$.

A separate exact optimization result shows that the classical Thue--Morse
phase $c=1/2$ minimizes the fourth-moment growth rate but is a strict local
maximum for the sixth-moment rate. We characterize the two global
sixth-moment minimizers by a degree-fourteen polynomial and certify their
uniqueness using rational Sturm and interval arithmetic. One minimizer is
approximately $0.364919195551255$.

The article supplies proofs, explicit matrices, recurrence formulas,
reproducible exact-arithmetic checks, a connection to the Fabius/Rvachev
Fourier program, and a proposed research agenda. The new claims are
research results argued here, not independently refereed conclusions or
Lean-verified theorems. Prior moment-matrix methods are explicitly credited.
\end{abstract}
\vfill
{\small\color{muted} Repository snapshot: \nolinkurl{37e61c1fdec28c6e7ab7ff445043993077b42cbd}.\\
Package: editable source, compiled PDF, two exact-arithmetic checkers, and their receipts.}
\end{titlepage}
% ed. (2026-09-30): page anchors back on after the title page.
\hypersetup{pageanchor=true}
\tableofcontents
\newpage

\section{The research question and the contribution}
\label{sec:intro}

\subsection{A concrete question in the literature}
The family of singular potentials
\[
 T(x)=2x\pmod 1,\qquad
 g_c(x)=\cos^2\!\pi(x-c),\qquad \psi_c=\log g_c
\]
is a natural deformation of the classical Thue--Morse sine product.
Gohlke, Kesseb\"ohmer, and Schindler establish real-analytic dependence on
$c$ for the pressure at order two and explicitly leave stronger regularity
at general positive orders open \cite[Theorem~2.7 and following discussion]{GKS}.
Their pressure--spectrum identity also connects the question to generalized
Thue--Morse diffraction measures \cite[Theorem~2.6]{GKS}.

Here the integer cases admit a particularly effective answer. Although
$\psi_c$ has a logarithmic singularity, $\exp(m\psi_c)=g_c^m$ is a
trigonometric polynomial when $m$ is a positive integer. The associated
transfer operator therefore has a small invariant space. The substantive
point is not merely writing down this space: a complete argument must
identify its spectral radius with pressure, prove simplicity at every phase,
and handle the exceptional phase $c=0$ without silently assuming
irreducibility there.

\subsection{Relationship to the repository}
The inspected \repo\ snapshot contains a large Fabius/Rvachev development
and a Fourier-decay synthesis \cite{ProveIt,RepoFourier}. The synthesis
already develops lacunary sine cocycles, pressure-based moment exponents,
and several distinct interpretations of Fourier decay. These are not
reintroduced as new results here. We instead study phase dependence and
phase optimization of the finite cocycle underlying that program.

Our arguments are self-contained at the level of finite-dimensional
spectral theory and elementary Fourier analysis. They do not assume that
a claim is correct merely because it appears in a repository draft. In
particular, none of the conclusions below requires a general
infinite-dimensional Ruelle--Perron--Frobenius theorem for a singular
potential. The measure-theoretic corollaries are explicitly separated and
uses the pressure--spectrum theorem of \cite{GKS}.

\subsection{What is old, what is proved here, and what remains open}
Finite recurrences for the classical even moments are established work.
Mauduit, Montgomery, and Rivat obtained an order-$m$ recurrence and an
$m\times m$ matrix for the $2m$th classical moment
\cite{MMR,MontgomeryTalk}. Bettin and Drappeau subsequently studied the
large-moment growth rate and answered a question from that work \cite{BD}.
The present $(2m-1)\times(2m-1)$ Fourier model treats arbitrary phases;
at the classical phase, reflection reduces the relevant space back to
dimension $m$.

\begin{center}
\begin{tabular}{@{}p{.30\textwidth}p{.62\textwidth}@{}}
\toprule
Status & Content \\
\midrule
Established background & Classical moment recurrences; the Fabius sinc product;
order-two phase analyticity; the cited measure--pressure identity.\\[2mm]
Proved in this article & All-positive-integer phase analyticity and a uniform
finite-size spectral expansion; the missing $2m$th Taylor coefficient;
the exact global sixth-moment phase optimum.\\[2mm]
Additional structural results & The entire spectrum at the atomic phase;
a phase-independent determinant; exact arbitrary-phase recurrences and
algebraicity statements.\\[2mm]
Not established here & Real-analyticity for every noninteger positive order;
large-$m$ optimal-phase asymptotics; independent priority or formal proof
assistant verification.\\
\bottomrule
\end{tabular}
\end{center}

The literature and repository search did not locate the all-integer
regularity theorem, the all-order missing-coefficient identity, or the
sixth-moment optimum in the forms proved below. This is a bounded novelty
assessment, not a guarantee that no equivalent result exists elsewhere.
The appropriate claim is a proposed research contribution with explicit
proofs and reproducible certificates, rather than an unqualified assertion
of a major breakthrough.

\section{Conventions and the main theorems}
\label{sec:main}
All logarithms are natural. We identify $\TT=\RR/\ZZ$ and put
$e(x)=\exp(2\pi\ii x)$. Write $\mathcal D_n$ for the $2^n$ dyadic
intervals of length $2^{-n}$, using closed intervals when taking suprema.
Shared endpoints do not affect the integrals.

For an integer $m\geq1$, define
\begin{align}
 F_{m,n,c}(x)&=\prod_{j=0}^{n-1}g_c(T^jx)^m,\\
 Z_{m,n}(c)&=\sum_{I\in\mathcal D_n}\sup_{x\in I}F_{m,n,c}(x),\\
 p_m(c)&=\lim_{n\to\infty}\frac1n\log Z_{m,n}(c),
 \label{eq:pressure-definition}
\end{align}
where existence of the last limit will be proved. This definition agrees
with positive-order topological pressure in the conventions of
\cite{GKS}. No use is made of the source's order-zero normalization:
with natural logarithms, the zero-potential pressure is $\log2$.

The phase-deformed finite product and its moments are
\begin{equation}
 Q_{n,c}(x)=\prod_{j=0}^{n-1}\bigl(1+e(2^jx-c)\bigr),\qquad
 M_{m,n}(c)=\int_0^1|Q_{n,c}(x)|^{2m}\dd x.
 \label{eq:Q}
\end{equation}
At $c=1/2$ this is the usual minus-sign Thue--Morse product. The elementary
normalization
\begin{equation}
 |Q_{n,c}(x)|^{2m}=4^{mn}F_{m,n,c}(x)
 \label{eq:normalization}
\end{equation}
will account for every power of two below.

\begin{theorem}[All-integer phase regularity and uniform moment asymptotics]
\label{thm:main-analytic}
For every integer $m\geq1$, there is an explicit matrix $A_m(c)$ of order
$2m-1$ with a positive, algebraically simple eigenvalue $\rho_m(c)$
strictly larger in modulus than every other eigenvalue. Moreover:
\begin{enumerate}[label=\textnormal{(\roman*)}]
\item $p_m(c)=\log\rho_m(c)$ exists and is real-analytic, even, and
one-periodic in $c$.
\item $\spec A_m(0)=\{2^{-j}:0\leq j\leq2m-2\}$ and
$\det A_m(c)=2^{-(m-1)(2m-1)}$ for every $c$.
\item There are a positive real-analytic function $b_m$ and constants
$K_m<\infty$, $0<\eta_m<1$ such that, uniformly in $c$ and $n\geq0$,
\[
 M_{m,n}(c)=b_m(c)\bigl(2^{2m-1}\rho_m(c)\bigr)^n
 \bigl(1+\varepsilon_{m,n}(c)\bigr),\qquad
 |\varepsilon_{m,n}(c)|\leq K_m\eta_m^n.
\]
\end{enumerate}
The constants are allowed to depend on $m$.
\end{theorem}

\begin{theorem}[The missing Taylor coefficient]
\label{thm:main-jet}
For $m\geq1$, let
\[
 R_m=\frac{2(2^{2m}-1)\zeta(2m)}{\pi^{2m}}.
\]
As $c\to0$, with $t=\pi c$,
\begin{align}
 \rho_m(c)&=\cos^{2m}t+R_m t^{2m}+O_m(t^{2m+2}),
 \label{eq:rho-main-jet}\\
 p_m(c)&=2m\log\cos t+R_m t^{2m}+O_m(t^{2m+2}).
 \label{eq:p-main-jet}
\end{align}
In particular,
\begin{equation}
 [c^{2m}]p_m(c)=0,
 \qquad
 p_m(c)=-2m\sum_{k=1}^{m-1}
 \frac{(2^{2k}-1)\zeta(2k)}{k}\,c^{2k}
 +O_m(c^{2m+2}).
 \label{eq:missing-main}
\end{equation}
For $m=1$, $p_1$ is identically zero.
\end{theorem}

\begin{theorem}[The fourth and sixth moments have different optimal phases]
\label{thm:main-optimum}
The unique minimizer of $\rho_2(c)$ modulo one is $c=1/2$, with
\[
 \rho_2(1/2)=\frac{1+\sqrt{17}}8.
\]
In contrast, $c=1/2$ is a strict local maximum of $\rho_3(c)$.
There are exactly two global minimizers of $\rho_3$ modulo one:
\[
 c_*=\frac{\arccos C_*}{2\pi},\qquad 1-c_*,
\]
where $C_*$ is the unique root in
$(-0.660930941723,-0.660930941722)$ of the explicit degree-fourteen
polynomial in \cref{eq:resultant-R}. Numerically,
\begin{align*}
 c_*&\approx0.364919195551255,\\
 \rho_3(c_*)&\approx0.430250310745843,\\
 32\rho_3(c_*)&\approx13.76800994386699.
\end{align*}
The two other phases $1/3$ and $1/2$ already give the exact separation
\begin{equation}
 \rho_3(1/3)<\frac7{16}<\rho_3(1/2).
 \label{eq:separator-main}
\end{equation}
The decimal for $C_*$ is certified by a rational interval; the additional
numerical displays are approximations, not the basis of the proof.
\end{theorem}

\section{The exact finite transfer model}
\label{sec:finite}

\subsection{An invariant Fourier space}
Define the transfer operator using a sum, rather than an average:
\begin{equation}
 (\Lop_{m,c}f)(x)=\sum_{Ty=x}g_c(y)^m f(y).
 \label{eq:transfer}
\end{equation}
Let
\[
 I_m=\{1-m,\ldots,m-1\},\qquad
 \mathcal E_m=\operatorname{span}_{\CC}\{e(rx):r\in I_m\}.
\]
The real subspace consists of those Fourier coefficients satisfying
$v_{-r}=\overline{v_r}$.

The binomial expansion gives
\begin{equation}
 g_c(x)^m=\sum_{j=-m}^m a_j e(jx-jc),\qquad
 a_j=4^{-m}\binom{2m}{m+j}.
 \label{eq:weight-coeff}
\end{equation}
We set $a_j=0$ outside $[-m,m]$. Since
\[
 \sum_{Ty=x}e(ky)=
 \begin{cases}2e(kx/2),&k\text{ even},\\0,&k\text{ odd},\end{cases}
\]
we obtain the following exact representation.

\begin{proposition}[Fourier matrix]
\label{prop:matrix}
The space $\mathcal E_m$ is invariant under $\Lop_{m,c}$. In the basis
indexed by $I_m$, its matrix is
\begin{equation}
 \boxed{\quad
 A_m(c)_{k,r}=2a_{2k-r}\,e(-(2k-r)c),\qquad k,r\in I_m.
 \quad}
 \label{eq:matrix}
\end{equation}
\end{proposition}
\begin{proof}
The product of a degree-$m$ trigonometric polynomial and one of degree at
most $m-1$ has degree at most $2m-1$. Summing over inverse branches
removes odd Fourier frequencies and halves the even ones. The resulting
degree is at most $m-1$. Applying this observation to a basis function
and using \cref{eq:weight-coeff} gives \cref{eq:matrix}.
\end{proof}

For example, with $z=e(c)$ and the order $-1,0,1$,
\begin{equation}
 A_2(c)=
 \begin{pmatrix}
 z/2&z^2/8&0\\
 1/(2z)&3/4&z/2\\
 0&1/(8z^2)&1/(2z)
 \end{pmatrix}.
 \label{eq:A2}
\end{equation}
The entries of this matrix need not be real or nonnegative. Positivity
below always means positivity of functions, not entrywise positivity of
the Fourier matrix.

\subsection{Exact moment recovery}
Let $v^{(0)}$ denote the coordinate vector of the constant function one,
and let $\ell_0$ extract the zeroth Fourier coefficient. Iterating the
transfer operator and changing variables gives
\begin{align}
 (\Lop_{m,c}^n1)(x)&=\sum_{T^ny=x}F_{m,n,c}(y),\\
 \ell_0 A_m(c)^n v^{(0)}
 &=2^n\int_0^1F_{m,n,c}(x)\dd x.
 \label{eq:integral-transfer}
\end{align}
Combining with \cref{eq:normalization} yields
\begin{equation}
 \boxed{\quad
 M_{m,n}(c)=2^{(2m-1)n}\ell_0 A_m(c)^n v^{(0)}.
 \quad}
 \label{eq:exact-moment}
\end{equation}
This identity holds for every $n\geq0$ and every phase, without an
asymptotic approximation. In particular, $m=1$ gives $A_1(c)=(1)$ and
$M_{1,n}(c)=2^n$, independently of $c$.

\section{Simplicity away from the atomic phase}
\label{sec:perron}

\subsection{The positive cone}
On the real subspace of $\mathcal E_m$, use the cone
\[
 \mathcal K_m=\{f\in\mathcal E_m:f(x)\geq0\text{ for all }x\in\TT\}.
\]
It is closed, pointed, generating, and has nonempty interior: the strictly
positive functions form its interior. For example, adding a sufficiently
large constant expresses any real trigonometric polynomial as a
difference of two members of the cone. The operator $\Lop_{m,c}$
preserves this cone.

We use the elementary finite-dimensional cone version of the Perron
principle: a linear map preserving a closed, pointed, solid cone has an
eigenvector in the cone at its spectral radius. One justification is to
add $\epsilon$ times a rank-one strongly positive map. Brouwer's fixed
point theorem on a compact normalized cone base gives a positive
eigenvector for the perturbed map; its eigenvalue dominates all others
by the order-unit norm. Letting $\epsilon\downarrow0$ and using continuity
of the finite spectrum gives the stated principle.

For the two inverse branches, the values of $g_c$ are $a$ and $1-a$,
where $0\leq a\leq1$. Therefore
\begin{equation}
 2^{1-m}\leq\Lop_{m,c}1=a^m+(1-a)^m\leq1.
 \label{eq:row-bound}
\end{equation}
Iteration and the spectral-radius formula imply
\begin{equation}
 2^{1-m}\leq\rho_m(c)\leq1.
 \label{eq:rho-bounds}
\end{equation}
In particular, the cone eigenvalue is positive.

\subsection{A one-zero backward-tree lemma}
The weight $g_c^m$ vanishes at exactly one point,
$s=c+1/2\pmod1$. This gives a short replacement for an inappropriate
entrywise-matrix irreducibility argument.

\begin{lemma}[Finite backward-invariant sets]
\label{lem:tree}
Suppose that $Z\subset\TT$ is nonempty and finite and
\[
 T^{-1}Z\setminus\{s\}\subseteq Z.
\]
Then $Z=\{0\}$ and $s=1/2$.
\end{lemma}
\begin{proof}
Writing $r=|Z|$, the left side has at least $2r-1$ points, so
$2r-1\leq r$ and $r=1$. If $Z=\{x\}$, one inverse image of $x$
must be $x$ itself, and the other must be the excluded point $s$.
Thus $Tx=x$. The only fixed point of doubling on $\TT$ is zero;
its two inverse images are $0$ and $1/2$.
\end{proof}

\begin{proposition}[Strictly positive Perron function]
\label{prop:h-positive}
If $c\neq0\pmod1$, every nonzero cone eigenfunction
$\Lop_{m,c}h=\rho_m(c)h$ is strictly positive.
\end{proposition}
\begin{proof}
A nonzero trigonometric polynomial has only finitely many zeros on the
circle. If $Z$ is the zero set of $h$, nonnegativity in the eigenvalue
equation implies $T^{-1}Z\setminus\{s\}\subseteq Z$. A nonempty
$Z$ would force $s=1/2$ by \cref{lem:tree}, hence $c=0$, a contradiction.
\end{proof}

\subsection{Excluding other peripheral eigenvalues}
\begin{proposition}[Simple dominant eigenvalue]
\label{prop:peripheral}
For $c\neq0$, the eigenvalue $\rho_m(c)$ is algebraically simple and
is the only eigenvalue of its modulus.
\end{proposition}
\begin{proof}
Let $h>0$ be a Perron eigenfunction and let $\Lop f=\lambda f$,
where $|\lambda|=\rho=\rho_m(c)$. Put
\[
 R=\max_{x\in\TT}\frac{|f(x)|}{h(x)}.
\]
If $f\neq0$, then $R>0$. The function
$H=R^2h^2-|f|^2$ is a nonnegative trigonometric polynomial. Whenever
$H(x)=0$, equality must hold in
\[
 |\lambda f(x)|
 \leq\sum_{Ty=x}g_c(y)^m|f(y)|
 \leq R\sum_{Ty=x}g_c(y)^mh(y)
 =R\rho h(x).
\]
Consequently, every inverse image of $x$ other than $s$ also lies in
the zero set of $H$. If $H$ were nonzero, its zero set would be finite
and nonempty, contradicting \cref{lem:tree}. Hence $|f/h|=R$
everywhere.

Write $F=f/h$. The normalized eigenvalue equation is a convex
combination of complex numbers of constant modulus. Equality in the
triangle inequality forces
\[
 F(y)=\frac\lambda\rho F(Ty)
 \quad\text{for }y\neq s.
\]
Continuity extends this identity to $s$. At the fixed point zero,
$F(0)\neq0$ gives $\lambda=\rho$. It then follows that $F=F\circ T$.
Every dyadic inverse image of zero has the same value as zero; these
points are dense, so $F$ is constant. The eigenspace is therefore
one-dimensional and no other peripheral eigenvalue exists.

Finally,
\[
 \|f\|_h=\max_x|f(x)|/h(x)
 \quad\Longrightarrow\quad
 \|\rho^{-1}\Lop f\|_h\leq\|f\|_h.
\]
A nontrivial Jordan block at $\rho$ would cause unbounded polynomial
growth of the powers of $\rho^{-1}\Lop$, which this estimate forbids.
Thus the eigenvalue is algebraically simple.
\end{proof}

\section{The atomic phase: spectrum and an exact Perron function}
\label{sec:atomic}
The backward-tree argument correctly identifies $c=0$ as exceptional.
Rather than treating this as an ignorable endpoint, we diagonalize its
spectrum and identify the positive Perron function explicitly.

\subsection{A binomial finite-difference diagonalization}
\begin{theorem}[Complete base-phase spectrum]
\label{thm:base-spectrum}
For every $m\geq1$,
\[
 \spec A_m(0)=\{1,2^{-1},2^{-2},\ldots,2^{-(2m-2)}\},
\]
and all eigenvalues are simple.
\end{theorem}
\begin{proof}
Let $J$ be a centered binomial random variable with
$\Pr(J=j)=a_j$, $-m\leq j\leq m$. Its two parity classes each have
probability $1/2$, as follows from $(1-1)^{2m}=0$.
For a function $P$ on $I_m$, \cref{eq:matrix} gives
\begin{equation}
 (A_m(0)^{\mathsf T}P)(r)
 =\mathbb E\!\left[P\left(\frac{r+J}{2}\right)
     \,\middle|\, J\equiv r\pmod2\right].
 \label{eq:binomial-markov}
\end{equation}
Every admissible $(r+J)/2$ belongs to $I_m$.

For every polynomial $U$ of degree less than $2m$,
\[
 \sum_{j=-m}^m(-1)^j a_jU(j)=0,
\]
because this is a constant multiple of a $2m$th finite difference.
Therefore the conditional moments of $J$ of degrees at most $2m-2$
are equal to the unconditional moments. If $P(r)=r^k$ for
$0\leq k\leq2m-2$, the right side of \cref{eq:binomial-markov} is a
polynomial in $r$ of degree $k$ with leading coefficient $2^{-k}$.

The evaluations of $1,r,\ldots,r^{2m-2}$ form a basis of functions on
$I_m$ by the Vandermonde determinant. In this basis the transpose
matrix is triangular with diagonal $2^{-k}$. The diagonal values are
distinct, proving the theorem.
\end{proof}

\subsection{The periodized-sinc eigenfunction}
We use $\sinc u=\sin u/u$, with removable value one at zero.

\begin{theorem}[The base-phase Perron function]
\label{thm:H}
The unique eigenfunction at eigenvalue one normalized by $H_m(0)=1$ is
\begin{equation}
 H_m(x)=\sum_{k\in\ZZ}\sinc^{2m}\!\bigl(\pi(x+k)\bigr).
 \label{eq:H-sinc}
\end{equation}
It belongs to $\mathcal E_m$ and is strictly positive on $\TT$.
Equivalently,
\begin{equation}
 H_m(x)=\frac1{(2m-1)!}
 \sum_{r=1-m}^{m-1}\eul{2m-1}{m-1+r}e(rx),
 \label{eq:H-eulerian}
\end{equation}
where the brackets denote Eulerian numbers.
\end{theorem}
\begin{proof}
The series in \cref{eq:H-sinc} converges uniformly on a period: its tails
are bounded by a constant times $\sum_{|k|\geq K}|k|^{-2m}$.
It is positive away from integers because every nonzero term is positive,
and equals one at integers.

For completeness, its finite Fourier expansion follows from the Fourier
transform of $\sinc^{2m}(\pi x)$. The transform of $\sinc(\pi x)$ is the
indicator of $[-1/2,1/2]$ under the $e^{-2\pi\ii x\xi}$ convention.
Thus the transform of the $2m$th power is the density $B_{2m}$ of the
sum of $2m$ independent uniforms on $[-1/2,1/2]$. One may justify the
product formula first in $L^2$ and then by convolution, since the even
power is integrable. The density has support $[-m,m]$ and vanishes
at its endpoints. Periodization therefore has no Fourier frequencies
outside $I_m$.

At an integer $r$ in that interval, the elementary inclusion--exclusion
formula for the volume of a clipped simplex gives
\[
 B_{2m}(r)=\frac1{(2m-1)!}
 \sum_{j=0}^{m+r}(-1)^j\binom{2m}{j}(m+r-j)^{2m-1}
 =\frac1{(2m-1)!}\eul{2m-1}{m-1+r}.
\]
The Fourier series of the periodization is absolutely finite, giving
\cref{eq:H-eulerian}. Alternatively, its coefficients follow by integrating
the uniformly convergent periodized series over one period.

It remains to verify the eigenvalue equation, which can be done without
Fourier transforms. For noninteger $x$,
\[
 H_m(x)=\frac{\sin^{2m}\pi x}{\pi^{2m}}
 \sum_{k\in\ZZ}(x+k)^{-2m}.
\]
In each term of $\Lop_{m,0}H_m$, use
$\sin\pi y\cos\pi y=\tfrac12\sin2\pi y$.
The branches $y=x/2$ and $y=(x+1)/2$ yield the even and odd terms of
the displayed sum after rescaling. Their sum is exactly $H_m(x)$.
Continuity gives the identity at integers. Uniqueness follows from
\cref{thm:base-spectrum}.
\end{proof}

Examples are
\[
 H_1=1,\qquad
 H_2(x)=\frac{2+\cos2\pi x}{3},\qquad
 H_3(x)=\frac{66+52\cos2\pi x+2\cos4\pi x}{120}.
\]
In particular,
\begin{equation}
 H_m(1/2)=\frac1{\pi^{2m}}\sum_{k\in\ZZ}(k+1/2)^{-2m}
 =\frac{2(2^{2m}-1)\zeta(2m)}{\pi^{2m}}=R_m.
 \label{eq:H-half}
\end{equation}
The sum of the coefficients in \cref{eq:H-eulerian} is one, as it must
be from the normalization at zero.

\section{Analyticity, pressure identification, and uniform errors}
\label{sec:analytic}

\subsection{Analytic perturbation in finite dimensions}
\begin{proposition}[Global analytic Perron data]
\label{prop:analytic-data}
For fixed $m$, the functions $\rho_m(c)$ and $p_m(c)=\log\rho_m(c)$
are real-analytic once pressure is identified below. There are also a
real-analytic rank-one spectral projector $P_m(c)$ and a real-analytic
positive eigenfunction $h_{m,c}$ normalized by $h_{m,c}(0)=1$.
\end{proposition}
\begin{proof}
Every entry of $A_m(c)$ is real-analytic in $c$.
\Cref{prop:peripheral,thm:base-spectrum} show that the positive dominant
eigenvalue is algebraically simple at every real phase. The implicit
function theorem applied to the characteristic polynomial therefore
gives a local analytic eigenvalue. The locally selected branches agree
on overlaps because the positive dominant eigenvalue is unique.
The rank-one projector can equivalently be obtained by integrating
$(zI-A_m(c))^{-1}$ around a small circle enclosing that eigenvalue.

The projector is positive as an operator on functions: the spectral gap
implies $\rho_m(c)^{-n}\Lop_{m,c}^n\to P_m(c)$, and each normalized
power is positive. If $h>0$ is any positive Perron function, then
$1\geq h/\|h\|_\infty$, so
$P_m(c)1\geq h/\|h\|_\infty>0$. Thus
\[
 h_{m,c}=\frac{P_m(c)1}{(P_m(c)1)(0)}
\]
is a well-defined positive analytic normalization, including at $c=0$
by \cref{thm:H}. Reflection $x\mapsto-x$ conjugates the phase $c$
operator to the phase $-c$ operator. Hence
\begin{equation}
 \rho_m(-c)=\rho_m(c),\qquad
 h_{m,-c}(x)=h_{m,c}(-x).
 \label{eq:reflection}
\end{equation}
Periodicity is immediate from the weights.
\end{proof}

\begin{proposition}[Uniform spectral estimate]
\label{prop:uniform}
For fixed $m$ there are constants $K<\infty$ and $0<\eta<1$ such that
\[
 \left\|\rho_m(c)^{-n}A_m(c)^n-P_m(c)\right\|
 \leq K\eta^n
\]
for every real $c$ and every $n\geq0$, in any fixed matrix norm.
\end{proposition}
\begin{proof}
Set $B(c)=A_m(c)/\rho_m(c)$ and $N(c)=B(c)-P_m(c)$.
Every eigenvalue of $N(c)$ has modulus strictly less than one.
Continuity of the finite spectrum and compactness of the phase circle
give $r_0=\max_c\rho(N(c))<1$. Choose $r_0<\eta<1$.
The resolvent $(zI-N(c))^{-1}$ is uniformly bounded for $|z|=\eta$
and all $c$, again by compactness. Cauchy's integral formula for matrix
powers gives $\|N(c)^n\|\leq K\eta^n$. Since $P_mN=NP_m=0$
and $P_m^2=P_m$, this gives the result for $n\geq1$; enlarge $K$ to
include $n=0$.
\end{proof}

Applying $\ell_0$ and $v^{(0)}$ to this estimate and using
\cref{eq:exact-moment} proves the moment expansion in
\cref{thm:main-analytic}, with
\begin{equation}
 b_m(c)=\ell_0P_m(c)v^{(0)}>0.
 \label{eq:b}
\end{equation}
The positivity proof in \cref{prop:analytic-data} and compactness show
that $b_m$ is bounded away from zero, so the absolute error may be
expressed as a uniform relative error.

At $c=0$, evaluation at zero is a left eigenfunctional because
$(\Lop_{m,0}f)(0)=f(0)$. Consequently
\begin{equation}
 P_m(0)f=f(0)H_m,\qquad
 b_m(0)=\int_0^1H_m(x)\dd x
 =\frac1{(2m-1)!}\eul{2m-1}{m-1}.
 \label{eq:bzero}
\end{equation}

\subsection{Why the finite eigenvalue is the actual pressure}
The following comparison avoids an unproved claim that pressure must
equal the spectral radius of an arbitrarily chosen invariant subspace.

\begin{lemma}[Cylinder suprema versus the integral]
\label{lem:variation}
If $F\geq0$ is a trigonometric polynomial of degree at most $D$ and
$N=2^n$, then
\begin{equation}
 N\int_0^1F\leq\sum_{I\in\mathcal D_n}\sup_I F
 \leq N\int_0^1F+\int_0^1|F'|
 \leq(N+4\pi D)\int_0^1F.
 \label{eq:variation-bound}
\end{equation}
\end{lemma}
\begin{proof}
On an interval $I$ of length $1/N$, the supremum is at least its average
and at most its average plus the variation of $F$ on $I$. Summing gives
the first two inequalities.

For the last inequality, factor $F(x)=|P(e(x))|^2$, with $P$ an ordinary
polynomial of degree at most $D$. This is the elementary
Fej\'er--Riesz factorization: reciprocal-conjugate roots of the Laurent
polynomial are paired, and roots on the unit circle have even
multiplicity because $F\geq0$; choosing one from each pair gives $P$.
Writing the same function as $P(x)=\sum_{j=0}^D b_je(jx)$,
Parseval and Cauchy--Schwarz give
\[
 \int|F'|\leq2\|P'\|_2\|P\|_2
 \leq4\pi D\|P\|_2^2=4\pi D\int F.
\]
\end{proof}

The polynomial $F_{m,n,c}$ has degree at most $m(2^n-1)$.
By \cref{eq:integral-transfer,eq:variation-bound},
\begin{equation}
 \ell_0A_m(c)^nv^{(0)}\leq Z_{m,n}(c)
 \leq(1+4\pi m)\ell_0A_m(c)^nv^{(0)}.
 \label{eq:pressure-sandwich}
\end{equation}
The positive leading coefficient \cref{eq:b} and the spectral estimate
now show that the logarithmic growth of both sides is $\log\rho_m(c)$.
This proves existence in \cref{eq:pressure-definition} and identifies
$p_m(c)$ with the finite eigenvalue. The proof of
\cref{thm:main-analytic}(i) and (iii) is complete.

\subsection{Strictness away from the atomic phase}
\begin{proposition}
\label{prop:strict}
For $m>1$, $\rho_m(c)=1$ if and only if $c=0\pmod1$.
\end{proposition}
\begin{proof}
We know $\rho_m(0)=1$. Suppose $\rho_m(c)=1$ and take a positive
Perron function $h$. At any maximizer $x$ of $h$, equality in
\[
 h(x)=\sum_{Ty=x}g_c(y)^mh(y)
 \leq\|h\|_\infty\sum_{Ty=x}g_c(y)^m
 \leq\|h\|_\infty
\]
forces $a^m+(1-a)^m=1$. For $m>1$ this requires $a\in\{0,1\}$.
The two inverse images are therefore $c$ and $c+1/2$, and $x=2c$.
The branch at $c$ has weight one, so $h(c)=\|h\|_\infty$ as well.
Every maximum must equal $2c$, hence $c=2c$ on the circle and $c=0$.
\end{proof}

\section{Algebraic structure: determinants and recurrences}
\label{sec:algebra}

\subsection{A phase-independent determinant}
Let $z=e(c)$ and write $A_m(z)$ for the Laurent-polynomial matrix.
From \cref{eq:matrix},
\[
 A_m(z)=\operatorname{diag}(z^{-2k})_{k\in I_m}\,
 A_m(1)\,\operatorname{diag}(z^r)_{r\in I_m}.
\]
Because $\sum_{k\in I_m}k=0$,
\begin{equation}
 \det A_m(z)=\det A_m(1)
 =\prod_{j=0}^{2m-2}2^{-j}
 =2^{-(m-1)(2m-1)}.
 \label{eq:det}
\end{equation}
This proves \cref{thm:main-analytic}(ii). In particular, none of the
finite-model eigenvalues can pass through zero as the phase changes.

\subsection{Characteristic polynomials in a single real parameter}
Put $d=2m-1$ and $C=\cos2\pi c$. The principal-minor expansion gives
\begin{equation}
 \chi_m(\lambda,z)=
 \sum_{S\subseteq I_m}(-1)^{|S|}\lambda^{d-|S|}
 \det A_m(1)[S,S]\,z^{-\sum_{k\in S}k}.
 \label{eq:principal-minors}
\end{equation}
Reflection $k\mapsto-k$ leaves the relevant determinants unchanged, so
this Laurent polynomial is invariant under $z\mapsto z^{-1}$.
It is therefore a polynomial in $C=(z+z^{-1})/2$ with rational
coefficients. The largest possible absolute subset sum in
\cref{eq:principal-minors} is $m(m-1)/2$, proving
\begin{equation}
 \chi_m(\lambda,C)\in\QQ[\lambda,C],\qquad
 \deg_C\chi_m\leq\frac{m(m-1)}2.
 \label{eq:degree-C}
\end{equation}
To distinguish phase and cosine parameters, set
\[
 \varrho_m(C)=\rho_m\!\left(\frac{\arccos C}{2\pi}\right).
\]
This is an analytic function of $C$ in a neighborhood of each point
of $[-1,1]$. Indeed, it is a simple root of
$\chi_m$ and the real implicit function theorem applies directly in
$(\lambda,C)$, including at the two endpoints.

\begin{corollary}[Algebraicity at arithmetic phases]
If $c=a/b$ is rational, then $\rho_m(c)$ and
$2^{2m-1}\rho_m(c)$ are algebraic numbers of degree at most
$(2m-1)\varphi(b)$ over $\QQ$.
More generally, they are algebraic whenever $\cos2\pi c$ is algebraic.
\end{corollary}
\begin{proof}
At a rational phase, $z=e(a/b)$ lies in a cyclotomic field of degree at
most $\varphi(b)$. The eigenvalue satisfies the degree-$2m-1$
characteristic polynomial over that field. The second assertion follows
from \cref{eq:degree-C}.
\end{proof}

\subsection{Exact linear recurrences}
Write
$\chi_m(\lambda,c)=\lambda^d+\gamma_1(c)\lambda^{d-1}
+\cdots+\gamma_d(c)$, and set
$u_n(c)=\ell_0A_m(c)^nv^{(0)}$.
Cayley--Hamilton yields
\begin{equation}
 u_{n+d}+\gamma_1u_{n+d-1}+\cdots+\gamma_du_n=0
 \qquad(n\geq0).
 \label{eq:recurrence}
\end{equation}
Since $M_{m,n}=2^{(2m-1)n}u_n$, this immediately gives a recurrence
for the actual moments as well. The stated order is an upper bound;
symmetry or cancellation can reduce the minimal order.

At $c=0$ and $c=1/2$, the weight is even. Starting from the constant
function, the operator stays in the even trigonometric subspace
$\operatorname{span}\{1,\cos2\pi x,\ldots,\cos2\pi(m-1)x\}$,
of dimension $m$. This explains compatibility with the earlier
classical order-$m$ recurrence \cite{MMR,MontgomeryTalk}; the larger
arbitrary-phase matrix is not claimed to improve that known bound.

\section{Delayed feedback and the missing coefficient}
\label{sec:jet}

\subsection{An exact equation at the fixed point}
Use the positive analytic normalization $h_{m,c}(0)=1$ from
\cref{prop:analytic-data}. Evaluating the transfer equation at zero
exposes its two inverse branches:
\begin{equation}
 \boxed{\quad
 \rho_m(c)=\cos^{2m}(\pi c)
   +\sin^{2m}(\pi c)\,h_{m,c}(1/2).
 \quad}
 \label{eq:feedback}
\end{equation}
The first summand comes from the fixed point zero. The second is the
feedback from the other inverse branch; its weight has a zero of order
$2m$ at the atomic phase.

Reflection \cref{eq:reflection}, together with $-1/2=1/2$ on the
circle, gives
\[
 h_{m,-c}(1/2)=h_{m,c}(1/2).
\]
Thus this value is an even analytic function of $c$. By
\cref{eq:H-half},
\begin{equation}
 h_{m,c}(1/2)=R_m+O_m(c^2).
 \label{eq:feedback-value}
\end{equation}
Substituting into \cref{eq:feedback}, and using
$\sin^{2m}t=t^{2m}+O_m(t^{2m+2})$, proves
\cref{eq:rho-main-jet}.

To pass to logarithms, factor out $\cos^{2m}t$ near zero. The quotient
is $1+R_mt^{2m}+O_m(t^{2m+2})$. Since its square-error term is
$O(t^{4m})\subseteq O(t^{2m+2})$ for every $m\geq1$, this proves
\cref{eq:p-main-jet}.

\subsection{Why the coefficient cancels exactly}
The Euler product for cosine is
\[
 \cos t=\prod_{j=0}^\infty
 \left(1-\frac{4t^2}{\pi^2(2j+1)^2}\right).
\]
It follows from the sine product by dividing $\sin2t/(2t)$ by
$\sin t/t$, or by pairing its zeros and using its value at zero.
For $|t|<\pi/2$, absolute convergence permits logarithmic expansion:
\begin{equation}
 \log\cos t
 =-\sum_{k\geq1}
 \frac{(2^{2k}-1)\zeta(2k)}{k\pi^{2k}}\,t^{2k}.
 \label{eq:logcos}
\end{equation}
The coefficient of $t^{2m}$ in $2m\log\cos t$ is
\[
 -\frac{2(2^{2m}-1)\zeta(2m)}{\pi^{2m}}=-R_m.
\]
This is precisely the negative of the feedback coefficient in
\cref{eq:p-main-jet}. The degree-$2m$ term vanishes, proving
\cref{thm:main-jet}.

This is not a claim that the whole pressure is flat to order $2m$.
For $m>1$, its leading curvature is
\begin{equation}
 p_m''(0)=-2m\pi^2<0.
 \label{eq:curvature}
\end{equation}
The cancellation removes one specific higher even coefficient, while
all lower even coefficients agree with the fixed-point contribution.

\subsection{Exact low-order expansions}
Let $t=\pi c$. The checker in the accompanying package computes the
following Taylor coefficients using exact rational arithmetic:
\begin{align}
 p_2(c)&=-2t^2+\frac{376}{45}t^6+O(t^8),
 \label{eq:jet2}\\
 p_3(c)&=-3t^2-\frac12t^4+\frac{213}{28}t^8+O(t^{10}),
 \label{eq:jet3}\\
 p_4(c)&=-4t^2-\frac23t^4-\frac8{45}t^6
 +\frac{16978624}{3075975}t^{10}+O(t^{12}).
 \label{eq:jet4}
\end{align}
The all-order theorem proves the missing terms $t^4,t^6,t^8$ in these
three displays. The next nonzero coefficients are finite exact
calculations rather than an asserted all-order pattern. For example,
$376/45$ in \cref{eq:jet2} is not inferred from the sign of a numerical
plot.

\begin{center}
\begin{tabular}{@{}cccccc@{}}
\toprule
$m$&2&3&4&5&6\\
\midrule
$R_m$&$1/3$&$2/15$&$17/315$&$62/2835$&$1382/155925$\\
Missing degree in $p_m$&4&6&8&10&12\\
\bottomrule
\end{tabular}
\end{center}

\section{Fourth- and sixth-moment phase optimization}
\label{sec:optimum}
It is enough to minimize $\rho_m(c)$ because the exponential moment
rate is $2^{2m-1}\rho_m(c)$, and the multiplier is phase-independent.
By evenness and periodicity, we may work in $C=\cos2\pi c\in[-1,1]$.
The characteristic-polynomial reductions used in this section are
checked directly from \cref{eq:matrix} in the supplied scripts.

\subsection{The fourth moment: a unique classical optimum}
For $m=2$,
\begin{equation}
 \chi_2(\lambda,C)=\lambda^3-(3/4+C)\lambda^2
 +(1/4+5C/8)\lambda-1/8.
 \label{eq:chi2}
\end{equation}
At every $C$,
\[
 \chi_2(5/8,C)=-9/512<0.
\]
Since the polynomial is positive for sufficiently large $\lambda$,
there is a positive root greater than $5/8$. The Perron root therefore
satisfies $\varrho_2(C)>5/8$.

At the Perron root, $\partial_\lambda\chi_2>0$. More generally,
for a real monic polynomial with a simple strictly dominant positive
root, every other real root contributes a positive difference, and
complex-conjugate roots contribute positive squared moduli. Hence
implicit differentiation gives
\[
 \frac{\dd\varrho_2}{\dd C}
 =-\frac{\varrho_2(5/8-\varrho_2)}{\partial_\lambda\chi_2(\varrho_2,C)}>0.
\]
Thus the minimum occurs uniquely at $C=-1$. At that phase,
\[
 \chi_2(\lambda,-1)
 =(\lambda+1/2)(\lambda^2-\lambda/4-1/4),
\]
and the positive root is $(1+\sqrt{17})/8$. This proves the
fourth-moment assertion in \cref{thm:main-optimum}.

\subsection{The sixth-moment characteristic polynomial}
For $m=3$, it is convenient to clear denominators and write
$f(\lambda,C)=2048\chi_3(\lambda,C)$:
\begin{align}
 f(\lambda,C)={}&2048\lambda^5
 -128(12C^2+15C+4)\lambda^4 \notag\\
 &+40(28C^3+24C^2+9C+1)\lambda^3 \notag\\
 &-(448C^3+261C^2-126C+37)\lambda^2 \notag\\
 &+(63C^2+21C-22)\lambda-2.
 \label{eq:f}
\end{align}
At the classical phase,
\begin{equation}
 f(\lambda,-1)=2(16\lambda^2+4\lambda-1)
 (64\lambda^3-20\lambda^2-6\lambda+1).
 \label{eq:half-factor}
\end{equation}
At phase $c=1/3$, where $C=-1/2$,
\begin{equation}
 4f(\lambda,-1/2)
 =8192\lambda^5+256\lambda^4-160\lambda^3
 -437\lambda^2-67\lambda-8=:P(\lambda).
 \label{eq:third-poly}
\end{equation}
There is exactly one positive root of $P$ by Descartes' rule of signs.
Since $P(0)<0$ and
\begin{equation}
 P(7/16)=405/64>0,
 \qquad
 (64\lambda^3-20\lambda^2-6\lambda+1)_{\lambda=7/16}
 =-3/32<0,
 \label{eq:separator}
\end{equation}
we get the exact strict separation
$\rho_3(1/3)<7/16<\rho_3(1/2)$.
This alone disproves the natural extrapolation that the classical phase
minimizes every positive even moment.

\subsection{The classical phase is a strict local maximum}
At $C=-1$, the Perron root is greater than $7/16$. The positive root
of the quadratic factor in \cref{eq:half-factor} is less than $1/4$,
so the Perron root belongs to the cubic factor. Reduction modulo that
cubic gives
\begin{equation}
 f_C(\rho,-1)=3\rho(720\rho^2-196\rho-41),
 \qquad \rho=\rho_3(1/2).
 \label{eq:fc-half}
\end{equation}
The quadratic in parentheses is positive for $\rho\geq7/16$:
its value at $7/16$ is $177/16$, and its derivative is positive there
and thereafter. Also $f_\lambda(\rho,-1)>0$ by the simple dominant
root property. Thus $\dd\varrho_3/\dd C<0$ at $C=-1$.

Writing $c=1/2+u$,
\[
 C=-1+2\pi^2u^2+O(u^4).
\]
It follows that $\rho_3(1/2+u)<\rho_3(1/2)$ for sufficiently small
nonzero $u$. The same is true of $p_3$ and of the sixth-moment
exponential growth factor.

\subsection{Certifying the global minimizer}
A stationary Perron root in $-1<C<1$ satisfies
\[
 f(\rho,C)=0,\qquad f_C(\rho,C)=0,
\]
because $f_\lambda(\rho,C)\neq0$. Eliminate $\rho$:
\begin{equation}
 \Res_\lambda(f,f_C)=-3057647616\,R(C),
 \label{eq:resultant-id}
\end{equation}
where
\begin{equation}
\begin{aligned}
 R(C)={}&478310547456C^{14}-263889031168C^{13}
 -2211948322304C^{12}\\
 &+2274923686784C^{11}+2772364425376C^{10}
 -3601461331552C^9\\
 &-2093687468208C^8+2695719201426C^7
 +1231497289965C^6\\
 &-1032213437124C^5-514151253800C^4
 +146152959398C^3\\
 &+104125846753C^2+11781400244C+375091954.
 \end{aligned}
\label{eq:resultant-R}
\end{equation}
A resultant alone is not enough: some of its zeros can represent
stationary non-Perron eigenvalues. The following certificate explicitly
removes those extraneous candidates.

Let $V(x)$ be the number of sign changes in the Sturm sequence of $R$
at $x$, ignoring zero entries. Exact rational arithmetic gives
$V(-1)=9$, $V(1)=5$, and the four isolating intervals below.
Every decimal endpoint in this table denotes an exact rational number.

\begin{center}
\small
\begin{tabular}{@{}p{.48\textwidth}cc@{}}
\toprule
Interval for $C$ & $V(\text{left}),V(\text{right})$
 & Common-root enclosure\\
\midrule
$(-0.660930941723,-0.660930941722)$ & $9,8$ & $(0.430,0.431)$\\
$(-0.070099130476,-0.070099130475)$ & $8,7$ & $(-0.152,-0.150)$\\
$(-0.068679451074,-0.068679451073)$ & $7,6$ & $(-0.134,-0.132)$\\
$(0.914028755782,0.914028755783)$ & $6,5$ & $(0.229,0.231)$\\
\bottomrule
\end{tabular}
\end{center}

Here is how the last column is certified, rather than numerically
assigned. The degree-one subresultant of $f$ and $f_C$ is
\begin{equation}
 169869312\bigl(A(C)\lambda+B(C)\bigr),
 \label{eq:linear-subresultant}
\end{equation}
where the integer polynomials $A,B$ are displayed in
\cref{sec:certificate-polynomials}. Subresultants are polynomial
combinations of the original two polynomials. The extracted content
$169869312$ is a nonzero constant, not a parameter-dependent factor.
Consequently every common root satisfies
$A(C)\lambda+B(C)=0$.

Rational interval Horner evaluation on each listed interval shows that
$A(C)$ is bounded away from zero and that $-B(C)/A(C)$ belongs to the
stated enclosure. Thus any stationary eigenvalue corresponding to the
last three rows is less than $1/4$. But \cref{eq:rho-bounds} gives
$\varrho_3(C)\geq1/4$ for all $C$, so none of those rows can represent a
stationary Perron eigenvalue.

A global Perron minimum exists by continuity. It lies in the interior
because the value at $C=-1/2$ is smaller than both endpoint values:
$\varrho_3(-1)>7/16$, $\varrho_3(1)=1$, and
$\varrho_3(-1/2)<7/16$. Every interior minimum is stationary. Hence the
first row must contain a stationary Perron point, and it is the only
one. This proves uniqueness of the minimizing $C_*$ and, by reflection,
exactly two minimizing phases modulo one.

For orientation, high-precision numerical evaluation gives
\[
 C_*\approx-0.6609309417223258423,\qquad
 \rho_3(1/2)\approx0.4443474454062066009.
\]
For comparison, the exponential sixth-moment rates are
\[
 32\rho_3(1/2)\approx14.21911825299861,\qquad
 32\rho_3(c_*)\approx13.76800994386699.
\]
The strict inequalities and uniqueness do not depend on these decimals.

\section{Consequences for dimensions and the Fabius program}
\label{sec:applications}

\subsection{Integer R\'enyi dimensions}
This subsection uses an external theorem, unlike the finite-operator
proofs above. Define the generalized Thue--Morse measure by the weak
limit of Riesz-product densities
\[
 \mu_c=\operatorname*{w-lim}_{n\to\infty}
 \left(\prod_{j=0}^{n-1}2g_c(T^jx)\right)\dd x.
\]
The existence and the identity
\begin{equation}
 \beta_{\mu_c}(m)=p_m(c)/\log2\qquad(c\neq0,\ m\geq1)
 \label{eq:imported-beta}
\end{equation}
are imported from \cite[Theorem~2.6]{GKS}. Here
\[
 \beta_\mu(q)=\lim_{n\to\infty}
 \frac1{n\log2}\log\sum_{I\in\mathcal D_n}\mu(I)^q,
\]
using a dyadic partition convention for the measure. For an integer
$m>1$, the R\'enyi dimension is
$D_m(\mu)=\beta_\mu(m)/(1-m)$.

\begin{corollary}[Analytic integer dimension functions]
\label{cor:dimensions}
With the cited measure--pressure identity, for every integer $m>1$,
\begin{equation}
 D_m(\mu_c)=\frac{p_m(c)}{(1-m)\log2}
 \label{eq:dimension}
\end{equation}
is real-analytic on the entire phase circle. Near $c=0$,
\[
 D_m(\mu_c)=\frac{2m}{(m-1)\log2}
 \sum_{k=1}^{m-1}\frac{(2^{2k}-1)\zeta(2k)}{k}c^{2k}
 +O_m(c^{2m+2}),
\]
and its $c^{2m}$ coefficient is zero.
\end{corollary}
\begin{proof}
For $c\neq0$ the formula is \cref{eq:imported-beta}. At $c=0$, the
measure is $\delta_0$ \cite{GKS}, so every positive-order dyadic mass
sum equals one and $D_m=0$. Also $p_m(0)=0$ for $m\geq1$ by our
base spectrum. Thus the equality extends at these integer orders.
Analyticity and the expansion follow from
\cref{thm:main-analytic,thm:main-jet}.
\end{proof}

The restriction to these orders matters: we have not asserted that the
measure--pressure identity extends at the atomic phase for all real
orders. Nor does the corollary give the entropy dimension at $m=1$ by
substituting into a formula with zero denominator.

\begin{corollary}[Distinct dimension-maximizing phases]
Under the same cited identity, $D_2(\mu_c)$ has its unique maximum
at $c=1/2$, whereas $D_3(\mu_c)$ has exactly two global maxima at
$c_*,1-c_*$. At $c=1/2$, $D_3$ has a strict local minimum.
\end{corollary}
\begin{proof}
The denominator in \cref{eq:dimension} is negative. Apply
\cref{thm:main-optimum} and the strict monotonicity of the logarithm.
\end{proof}

Thus a phase optimized for one integer spectral dimension need not be
optimized for the next. The sixth-moment computation is not merely an
alternative decimal evaluation of the known fourth-moment exponent.

\subsection{What the result says about the up-function}
Under the Fourier convention
$\widehat f(\xi)=\int f(x)e^{-2\pi\ii x\xi}\dd x$, the Rvachev
up-function has the classical transform
\begin{equation}
 \Phi(\xi)=\widehat{\operatorname{up}}(\xi)
 =\prod_{h=0}^{\infty}\sinc\!\left(\frac{\pi\xi}{2^h}\right).
 \label{eq:Phi}
\end{equation}
See \cite{Arias,RepoFourier}. Factoring its first $n$ dilated terms gives
for $x\neq0$
\begin{equation}
 \Phi(2^nx)=\Phi(x)\,(\pi x)^{-n}2^{-n(n+1)/2}
 \prod_{j=1}^n\sin(\pi2^jx).
 \label{eq:Phi-factor}
\end{equation}
The oscillatory factor is exactly
\[
 2^n\left|\prod_{j=1}^n\sin(\pi2^jx)\right|
 =|Q_{n,1/2}(2x)|.
\]
Since doubling preserves Lebesgue measure on the circle,
\begin{equation}
 \int_0^1\prod_{j=1}^n|2\sin(\pi2^jx)|^{2m}\dd x
 =M_{m,n}(1/2).
 \label{eq:fabius-moment}
\end{equation}
Thus the finite model directly computes the exponential rates of all
even moments of the stripped Fourier numerator.

The new phase theorems provide a controlled deformation of that
numerator and exact comparison phases. They can be used as analytic
benchmarks for pressure code, as finite algebraic test cases for
formalization, and as obstructions to assuming the classical phase is
universally moment-optimal.

There is an important limitation. The factors $x^{-n}$ and $\Phi(x)$ in
\cref{eq:Phi-factor} are not constant on a shell. Therefore
\cref{eq:fabius-moment} is not automatically an asymptotic for an
unweighted Fourier-shell integral. No new universal pointwise decay
exponent, annular maximum formula, or noninteger shell-pressure theorem
is claimed. The distinctions between these quantities are already
central to the repository's Fourier-decay draft \cite{RepoFourier}.

\section{Reproducible computation and a formalization route}
\label{sec:verification}

\subsection{Exact Taylor recursion}
The matrices give a rational algorithm for arbitrary finite jets at
$c=0$. Use $t=\pi c$ and write
\[
 A(t)=\sum_{n\geq0}A_nt^n,\quad
 v(t)=\sum_{n\geq0}v_nt^n,\quad
 \rho(t)=\sum_{n\geq0}\lambda_nt^n,
\]
where $\lambda_0=1$. By \cref{eq:matrix},
\begin{equation}
 (A_n)_{k,r}=(A_0)_{k,r}\frac{[-2\ii(2k-r)]^n}{n!}.
 \label{eq:An}
\end{equation}
Let $u$ be the row of ones, so $uv$ is evaluation at zero. Normalize
$uv_0=1$ and $uv_n=0$ for $n\geq1$. The vector $v_0$ is the
Eulerian coefficient vector in \cref{eq:H-eulerian}.

For each $n\geq1$, solve the same bordered linear system
\begin{equation}
 \begin{pmatrix}A_0-I&-v_0\\u&0\end{pmatrix}
 \begin{pmatrix}v_n\\\lambda_n\end{pmatrix}
 =\begin{pmatrix}
 -\sum_{k=1}^nA_kv_{n-k}
 +\sum_{k=1}^{n-1}\lambda_kv_{n-k}\\0
 \end{pmatrix}.
 \label{eq:bordered}
\end{equation}
The matrix is invertible because eigenvalue one is simple and
$uv_0=1$. Its entries lie in $\QQ(\ii)$, so every coefficient is exact.
If $p(t)=\log\rho(t)=\sum_{n\geq1}p_nt^n$, then differentiating
$\rho=\exp p$ gives the recursion
\begin{equation}
 p_n=\lambda_n-\sum_{k=1}^{n-1}\frac{k}{n}p_k\lambda_{n-k}.
 \label{eq:log-recursion}
\end{equation}
This is a finite algebraic way to check any desired coefficient without
numerically differentiating a nearly degenerate matrix.

\subsection{What was actually executed}
The package includes two Python checkers requiring SymPy. The execution
receipts were generated with SymPy~1.14.0. Neither checker uses floating
point arithmetic for its assertions.

\paragraph{\code{verify_pressure.py}.}
The default run passes 23 groups of checks. It verifies the base spectrum,
Eulerian Perron vector, and determinant for $1\leq m\leq10$; symbolic
phase-independent determinants and reflection for $1\leq m\leq4$;
exact Parseval moment identities for $1\leq m\leq4$, $0\leq n\leq6$,
and both signs; and exact Taylor coefficients through degree $2m+2$
for $2\leq m\leq6$. The receipt is \code{verification.json}.

\paragraph{\code{verify_minimum.py}.}
This checks the sixth-moment characteristic polynomial directly from
the matrix, computes the exact resultant and linear subresultant,
constructs a rational Sturm sequence, proves the four isolating intervals,
and bounds each candidate common eigenvalue by rational interval Horner
arithmetic. It also checks the separator values at $7/16$.
The receipt is \code{minimum_certificate.json}.

These finite checks are safeguards against transcription, scaling, and
low-order algebra mistakes. They do not replace the all-$m$ proofs of
analyticity or coefficient cancellation. The global sixth-moment theorem
does use the finite exact certificate as part of its proof; the integer
polynomials and verification recipe are supplied so it can be audited
independently.

\subsection{A practical Lean decomposition}
No Lean build was performed for this article. A proposed implementation
should keep the analytic and finite algebraic components separate.

The first layer would define finite Laurent polynomials and prove the
coefficient formula \cref{eq:matrix}, the moment identity, the
phase-independent determinant, and the characteristic polynomials for
$m=2,3$. These are finite algebraic targets and do not need a general
thermodynamic-formalism library.

A second layer would formalize the finite backward-tree lemma and the
fact that a nonzero trigonometric polynomial has finitely many zeros.
The positivity and peripheral-spectrum proof can then be phrased in
terms of functions and a weighted supremum norm. This avoids trying to
apply an entrywise-positive-matrix theorem to a complex Fourier matrix.

A third layer would prove the binomial finite-difference spectrum at
$c=0$ and the normalized $H_m$ formula. The finite Eulerian version may
be the easier route for exact values, whereas periodization makes the
analytic proof of $R_m$ conceptually transparent.

Finally, a spectral perturbation layer would deliver real-analytic
Perron data and feed the fixed-point equation \cref{eq:feedback} into a
power-series argument. The sixth-moment certificate could be imported as
explicit polynomial identities and rational sign checks, but a kernel
checked Sturm theorem would be preferable to trusting an external
computer-algebra receipt.

This order produces useful intermediate theorems even before the entire
article is formalized. It also preserves a visible boundary between
proof-assistant results, exact external computations, and proposed
analytic extensions.

\section{Further research questions}
\label{sec:questions}
The following questions are proposed continuations, not claims that they
have all been identified as open in the published literature. The first
is the remaining part of the explicit regularity question in \cite{GKS}.

\subsection{Noninteger regularity and the atomic threshold}
\begin{question}
For a fixed noninteger $q>0$, what is the exact regularity of
$c\mapsto P(q\psi_c)$, especially at $c=0$?
\end{question}
The integer proof relies on a finite Fourier cutoff that disappears for
noninteger $q$. The fixed-point equation suggests examining whether a
term of order $|c|^{2q}$ appears and whether special cancellations can
occur. This is a heuristic direction, not an established noninteger
expansion. A rigorous route would need a suitable eigenfunction space,
control of the singular branch, and a proof that the relevant eigenvalue
still represents pressure. Different regimes near $q=1/2$ deserve
separate treatment rather than extrapolating the integer argument.

\subsection{The first coefficient after the cancellation}
\begin{question}
Is $[t^{2m+2}]p_m(t/\pi)$ positive for every integer $m\geq2$?
Can it be written in a closed form in $m$?
\end{question}
The executed checks are positive for $2\leq m\leq6$.
Differentiating \cref{eq:feedback} two additional times reduces the
problem to the second phase derivative of $h_{m,c}(1/2)$ at zero.
The completely explicit base spectrum might permit a finite resolvent
sum with a useful sign structure.

\subsection{Optimal phases for higher moments}
\begin{question}
For each $m\geq4$, how many phases globally minimize $\rho_m(c)$?
How do those phases depend on $m$?
\end{question}
The change between $m=2$ and $m=3$ rules out a simple universal
classical-phase principle. The characteristic polynomial in $C$ gives
an exact elimination procedure at each fixed $m$, but resultant size
grows rapidly. A combination of cone bounds, symmetry, and certified
interval eigenvalue methods may scale better than full elimination.

\subsection{The large-moment phase landscape}
\begin{question}
Do minimizing phases converge as $m\to\infty$? Can their accumulation
points be characterized by ergodic optimization of $\log g_c$?
\end{question}
Existing large-moment results at the classical phase \cite{BD} should
be treated as input and comparison, not rediscovered. A new result
would have to control the phase uniformly. In particular, a fixed-$m$
spectral gap cannot simply be interchanged with the limit $m\to\infty$.

\subsection{Minimal recurrences and algebraic degree}
\begin{question}
What is the generic minimal recurrence order of $M_{m,n}(c)$, and
which phases cause it to drop below $2m-1$?
\end{question}
Reflection explains the reduction at $0$ and $1/2$, but other
cancellations could occur between the initial vector, spectral
projectors, and the integration functional. The appropriate finite
criterion is the rank of the corresponding Hankel or Krylov matrix,
not merely the degree of the characteristic polynomial.

\subsection{Quantitative gaps and certified errors}
\begin{question}
Can the constants $K_m,\eta_m$ in the uniform moment estimate be bounded
explicitly and efficiently in $m$?
\end{question}
Compactness proves their existence but is not a numerical certificate.
The one-zero argument proves strictness without a ready-made lower
bound on separation. An effective proof could use a bounded iterate
that contracts the positive cone, with a separate treatment near the
atomic phase using its exact spectrum.

\subsection{Beyond one zero and beyond base two}
\begin{question}
Which nonnegative trigonometric weights for $x\mapsto bx$ admit an
analogous all-phase simple finite Perron model?
\end{question}
The degree cutoff generalizes immediately. The subtle part is spectral
simplicity. With $s$ forbidden inverse images and a finite equality set
of size $r$, the counting argument becomes $br-s\leq r$. Its possible
small invariant sets should determine exceptional parameter loci.
This suggests a classification theorem organized by finite backward
orbits of the zero set.
% ed. (2026-09-30): a later package bears on this question
\begin{ednote}
\path{../Periodic_Anchors_Exact_Mixed_Moment_Phase_Diagrams/} (batch~68 of
the repository's \path{docs/incoming/}) computes exactly the pressure of
products of shifted digital masks whose phases fill complete periodic
orbits of $x\mapsto bx\bmod1$, with a common exponent on each orbit,
together with every equilibrium measure and the resulting mixed-moment
exponents, in every base. Such a product has several zeros, all mapped by $x\mapsto bx$ onto the selected orbits, and its pressure is a maximum of explicit orbit branches; no finite Perron model or spectral simplicity is proved. It bears on this question but does not
answer it: the phases are fixed, a single mask whose phase has period at
least two is not covered, and no regularity in the phase is proved.
\end{ednote}

\subsection{Phase sequences and matrix products}
\begin{question}
What changes when the phase is allowed to vary with scale, replacing
$c$ by $c_j$ in the $j$th factor?
\end{question}
The exact moment representation becomes a product of the same finite
matrices. Periodic phase sequences are finite algebraic models; more
general sequences lead to Lyapunov exponents and joint spectral radii.
This extension may distinguish optimization by a constant phase from
optimization by a scale-dependent protocol.

\subsection{A kernel-checked optimization certificate}
\begin{question}
Can the sixth-moment optimum be formalized with a small, transparent
certificate whose correctness is checked entirely in Lean?
\end{question}
The proof naturally splits into matrix identities, the analytic simple
root theorem, and a degree-fourteen real-root certificate. Establishing
this example would create a reusable workflow for moment optimization
problems elsewhere in \repo.

\section{Conclusions and scope}
The integer powers of the generalized Thue--Morse weight turn a singular
pressure problem into finite Fourier algebra. The key spectral step is
a one-zero backward-tree argument, supplemented by an exact treatment
of the atomic phase. Together they establish real-analytic phase
dependence for every positive integer order and a uniform exponential
moment expansion.

The fixed-point eigenvalue equation then explains a more unexpected
identity: feedback from the second branch enters precisely at degree
$2m$ and cancels that coefficient of the pressure. The sixth-moment
calculation shows that the phase geometry is genuinely richer than the
fourth-moment case: the classical phase ceases to minimize the rate,
and the true optimum has an explicit algebraic characterization.

These conclusions concern a concrete slice of the repository's
Fabius/Thue--Morse program. They do not resolve all positive real pressure
orders, do not determine a universal Fourier-decay law, and have not
been checked by a proof assistant or independent referee. The supplied
proofs and exact certificates make those remaining validation tasks
specific rather than leaving the central assertions at the level of
numerical evidence.

\appendix
\section{The sixth-moment certificate polynomials}
\label{sec:certificate-polynomials}
The polynomials in \cref{eq:linear-subresultant} are
\begin{align}
 A(C)={}&-34165039104C^{13}-10163951616C^{12}
 +145543077120C^{11}\notag\\
 &-39584395200C^{10}-236661608848C^9
 +88141321128C^8\notag\\
 &+214754071532C^7-68497376017C^6
 -115302016540C^5\notag\\
 &+21790610488C^4+33317271568C^3
 -2056400573C^2\notag\\
 &-4362939336C-301154570,
 \label{eq:A-certificate}\\[2mm]
 B(C)={}&3253813248C^{11}-2632283136C^{10}
 -7946511104C^9\notag\\
 &+9092577984C^8+8363143264C^7
 -10354590704C^6\notag\\
 &-4894201844C^5+5242980644C^4
 +2297293628C^3\notag\\
 &-801503452C^2-418577016C-27190504.
 \label{eq:B-certificate}
\end{align}
A reader can independently reproduce the certificate as follows.
Compute the resultant and polynomial subresultant sequence of
\cref{eq:f} and its $C$ derivative, eliminating $\lambda$.
Check \cref{eq:resultant-id,eq:linear-subresultant}. Construct the Sturm
sequence of $R$ using ordinary Euclidean remainders with reversed sign.
Evaluate its signs at $-1,1$ and the eight rational endpoints in the
table in \cref{sec:optimum}. Finally, interval-evaluate $A$ and $B$ by
Horner's rule, check that the denominator interval excludes zero, and
divide the interval for $-B$ by that for $A$.

Every operation in this recipe is addition, multiplication, comparison,
or exact division of rational numbers. Floating-point root finding,
plot sampling, and heuristic root matching are unnecessary.

\section{Build and verification instructions}
The source uses a standard \LaTeX{} installation with the Libertinus,
AMS, microtype, hyperref, and cleveref packages. No repository checkout
or network access is needed to compile it. From the package directory:
\begin{verbatim}
latexmk -pdf -interaction=nonstopmode -halt-on-error article.tex
python verify_pressure.py
python verify_minimum.py
\end{verbatim}
The last two commands require Python~3.10 or later and SymPy; the tested
version is recorded in \code{requirements.txt}. The algebraic checks
write JSON receipts to the working directory. Avoid running Python with
\code{-O}, because assertions are part of these verification scripts.
The article is self-contained; its bibliography URLs are for provenance
and further reading, not build dependencies.

\clearpage
\phantomsection
\addcontentsline{toc}{section}{References}
\begin{thebibliography}{9}
\bibitem{GKS}
P. Gohlke, M. Kesseb\"ohmer, and T. I. Schindler,
\emph{Generalized Thue--Morse measures: spectral and fractal analysis},
arXiv:2509.22109v1, 26 September 2025.
\url{https://arxiv.org/abs/2509.22109}.
Theorem~2.7 and its following paragraph give the regularity question;
Theorems~2.3 and 2.6 give the pressure identifications used for context
and the measure corollary. Version checked on 28 September 2026.

\bibitem{MMR}
C. Mauduit, H. L. Montgomery, and J. Rivat,
\emph{Moments of a Thue--Morse generating function},
Journal d'Analyse Math\'ematique \textbf{135} (2018), 713--724.
\url{https://doi.org/10.1007/s11854-018-0050-y}.
The classical moment-matrix result is also described in the primary
author presentation \cite{MontgomeryTalk}.

\bibitem{MontgomeryTalk}
H. L. Montgomery,
\emph{Moments of a Thue--Morse generating function},
CIRM research lecture recorded 12 February 2014.
\url{https://doi.org/10.24350/CIRM.V.18610603}.
\url{https://www.carmin.tv/en/collections/prime-numbers-new-perspectives-nombres-premiers-nouvelles-perspectives/video/moments-of-a-thue-morse-generating-function}.

\bibitem{BD}
S. Bettin and S. Drappeau,
\emph{Two arithmetic applications of perturbations of composition operators},
arXiv:2002.11417.
\url{https://arxiv.org/abs/2002.11417}.
Published version: \url{https://doi.org/10.1007/s11854-021-0184-1}.

\bibitem{Arias}
J. Arias de Reyna,
\emph{An infinitely differentiable function with compact support:
Definition and properties}, arXiv:1702.05442.
\url{https://arxiv.org/abs/1702.05442}.

\bibitem{ProveIt}
V. Reshetnikov and contributors, \emph{ProveIt}. Repository snapshot
\nolinkurl{37e61c1fdec28c6e7ab7ff445043993077b42cbd};
inspected 28 September 2026.
\url{https://github.com/VladimirReshetnikov/ProveIt/tree/37e61c1fdec28c6e7ab7ff445043993077b42cbd/Analysis/FabiusFunction}.

\bibitem{RepoFourier}
\emph{Fourier Decay of Rvachev's Up-Function: Exact products, pressure
spectra, Ces\`aro gauges, fluctuations, peaks, and valleys},
research synthesis in the \repo\ snapshot of \cite{ProveIt}.
Repository path:
\path{Analysis/FabiusFunction/docs/semi-formalized-research-frontiers/drafts/rvachev_up_fourier_decay/Rvachev_Up_Fourier_Decay/Rvachev_Up_Fourier_Decay.tex}.
The inspected source is dated 31 August 2026. It is used as repository
context, not as an independently refereed source or an unconditional
black-box proof of the present results.
\end{thebibliography}
\end{document}
